\documentclass{article}

\usepackage{amsmath}
\usepackage{tabularx}
\usepackage{booktabs}
\usepackage{hyperref}
\usepackage{enumitem}

\title{User Manual --- \progname}

\author{\authname}

\date{}

\input{../../Comments.text}
\input{../../Common.text}

\begin{document}

\maketitle

\begin{table}[hp]
\caption{Revision History} \label{TblRevisionHistory}
\begin{tabularx}{\textwidth}{llX}
\toprule
\textbf{Date} & \textbf{Developer(s)} & \textbf{Change}\\
\midrule
Apr 20, 2026 & Mirhoseininejad & Initial User Manual covering system
requirements, installation (backend and frontend), operation, usage
walk-through, troubleshooting, and uninstallation for \progname{}
v1.0.\\
\bottomrule
\end{tabularx}
\end{table}

\newpage

\tableofcontents

\newpage

\section{Introduction}

\progname{} is a browser-based exercise form-feedback tool. It uses
the user's webcam to estimate body pose, computes joint angles in
real time, and overlays a skeleton plus a rep counter and form-quality
indicator on the live video. This manual describes how to install and
use the reference implementation. It is written for a user who has
basic familiarity with a terminal but is not expected to have prior
Python or Node.js experience.

\section{System Requirements}

\subsection{Hardware}

\begin{itemize}
  \item A computer with an integrated or external USB webcam.
  \item At least 4~GB of free RAM.
  \item A recent CPU: Apple~Silicon (M1 or newer) or any x86-64 CPU
    from the last five years is sufficient.
  \item A display resolution of at least $1280\times720$.
\end{itemize}

\subsection{Software}

\begin{itemize}
  \item \textbf{Operating system:} macOS~12 or newer, or a modern
    Linux distribution (Ubuntu~22.04 or equivalent). Windows is
    expected to work but has not been exercised during CAS~741
    testing.
  \item \textbf{Python~3.9} (for the backend). Other Python~3 versions
    in the 3.9--3.11 range are likely to work but 3.9.6 is the
    reference version used during development.
  \item \textbf{Node.js~18} or newer (for the frontend; Vite requires
    a modern Node).
  \item \textbf{Redis} (optional, used by Django Channels for the
    WebSocket transport; the default configuration uses an in-memory
    channel layer that does not require Redis for single-process
    local use).
  \item \textbf{A modern browser} supporting the
    \texttt{getUserMedia} API --- Chrome, Edge, Firefox, or Safari in
    their current versions.
\end{itemize}

\section{Installation}

\subsection{Obtain the source}

Clone the repository and check out the submodule that contains the
application code:

\begin{verbatim}
git clone --recursive https://github.com/mmdmirh/cas741
cd cas741/FitCoachAR
\end{verbatim}

If the repository was cloned without the \texttt{--recursive} flag,
initialise the submodule explicitly:

\begin{verbatim}
git submodule update --init --recursive
\end{verbatim}

\subsection{Backend}

From the \texttt{FitCoachAR/src/backend} directory:

\begin{verbatim}
cd src/backend
python3 -m venv venv
source venv/bin/activate
pip install --upgrade pip
pip install -r requirements.txt
\end{verbatim}

The pinned dependencies include Django, Channels, MediaPipe (via the
\texttt{mediapipe-silicon} wheel on Apple~Silicon;
\texttt{mediapipe} on x86-64), NumPy, SciPy, and \texttt{pykalman}.

Apply the Django migrations:

\begin{verbatim}
python manage.py migrate
\end{verbatim}

\subsection{Frontend}

From the \texttt{FitCoachAR/src/frontend} directory, in a separate
terminal:

\begin{verbatim}
cd src/frontend
npm install
\end{verbatim}

\section{Running \progname}

\subsection{Start the backend}

From the \texttt{src/backend} directory, with the virtual environment
active:

\begin{verbatim}
python manage.py runserver
\end{verbatim}

The backend listens on \texttt{http://localhost:8000}. The WebSocket
endpoint used by the frontend is served on the same host.

\subsection{Start the frontend}

From the \texttt{src/frontend} directory:

\begin{verbatim}
npm run dev
\end{verbatim}

Vite prints the development URL (typically
\texttt{http://localhost:5173}). Open that URL in a supported browser.

\subsection{Grant camera permission}

The first time the application is opened, the browser asks for
permission to use the webcam. Camera permission is required;
\progname{} cannot function without it. If the permission was denied
previously, re-enable it through the browser's site-settings panel
and reload the page.

\section{Using the Application}

\subsection{Select an exercise}

The sidebar panel lists the supported exercises (squat, bicep curl).
Choose one before starting a session. The selection is sent to the
backend as the session parameter and determines which validity rules
the state machine applies.

\subsection{Position the camera}

For best results:

\begin{itemize}
  \item Place the camera so the full body is in frame for squats, or
    so the relevant arm is clearly visible for bicep curls.
  \item Stand 2--3 metres back from the camera.
  \item Use a plain background and well-diffused lighting.
  \item Wear form-fitting clothing where practical; very loose
    clothing reduces landmark tracking accuracy.
\end{itemize}

\subsection{Interpret the overlay}

The live view shows three elements:

\begin{itemize}
  \item \textbf{Skeleton overlay.} Lines drawn between the tracked
    body landmarks. A visible overlay means the pose backend is
    tracking you successfully. A flickering overlay indicates partial
    occlusion or poor lighting.
  \item \textbf{Rep counter.} Integer count of valid repetitions
    completed so far in the current session. Increments only when
    both the range-of-motion and timing constraints are satisfied.
  \item \textbf{Form-quality badge.} Short text feedback such as
    \textit{``Go deeper''} (partial range of motion) or
    \textit{``Too fast''} (rep completed in less than
    $t_{\min}$).
\end{itemize}

\subsection{Ending a session}

Close the browser tab or click the \textit{End session} control in
the sidebar. The WebSocket closes cleanly; no explicit logout is
required.

\section{Troubleshooting}

\begin{table}[h]
\centering
\begin{tabularx}{\textwidth}{p{5cm} X}
\toprule
\textbf{Symptom} & \textbf{Likely cause and resolution}\\
\midrule
Browser says ``camera in use by another application.'' &
  Close any other tab or application that is using the webcam and
  reload the page.\\
No skeleton appears on the video. &
  The pose backend did not initialise. Check the backend console
  for a MediaPipe import error. On Apple Silicon, verify
  \texttt{mediapipe-silicon} is installed rather than plain
  \texttt{mediapipe}.\\
Rep counter never increments. &
  Confirm you are reaching the full range of motion for the selected
  exercise (for squats, the knee must cross
  $\theta_{\text{low}} = 90^\circ$). Verify the camera sees your full
  lower body.\\
``Too fast'' shown on every rep. &
  The minimum rep duration $t_{\min}$ is 0.5~s by default. Slow down,
  or adjust $t_{\min}$ in the calibration parameters.\\
Frontend cannot reach the backend WebSocket. &
  Verify the backend is running on
  \texttt{localhost:8000} and that no local firewall is blocking
  loopback traffic.\\
The page loads but latency is very high ($<10$~FPS). &
  Check CPU load on the machine and close other heavy applications.
  See the Performance Report for the expected latency budget per
  backend.\\
\bottomrule
\end{tabularx}
\caption{Common problems and resolutions.}
\label{tab:trouble}
\end{table}

\section{Uninstallation}

To remove \progname{} from the machine:

\begin{enumerate}
  \item Stop the backend and frontend processes
    (\texttt{Ctrl+C} in each terminal).
  \item Deactivate and delete the backend virtual environment:
    \begin{verbatim}
deactivate
rm -rf src/backend/venv
    \end{verbatim}
  \item Delete the frontend dependency cache:
    \begin{verbatim}
rm -rf src/frontend/node_modules
    \end{verbatim}
  \item Delete the cloned repository itself if no longer needed:
    \begin{verbatim}
cd ../../../..
rm -rf cas741
    \end{verbatim}
\end{enumerate}

No system-level packages are installed by \progname{}, so no
system-wide cleanup is required.

\section{Getting Help}

Bug reports and questions can be filed as issues on the project
repository at
\href{https://github.com/mmdmirh/cas741}{github.com/mmdmirh/cas741}.
The supporting documentation --- Problem Statement, SRS, VnV Plan,
VnV Report, Module Guide, and Module Interface Specification --- is
in the \texttt{docs/} folder of the same repository.

\end{document}
